This chapter introduces the concepts on which the proposed system is built:
transformer-based Large Language Models and the prompting techniques used to
control them, the psychology of musical emotion that grounds the
emotion-to-key mapping, and the symbolic music representations used to encode
and export generated melodies.

\section{Large Language Models}
\label{sec:bg-llm}

Large Language Models are neural networks based on the Transformer
architecture \cite{vaswani17attention}, trained on internet-scale text
corpora with a self-supervised next-token prediction objective. The
Transformer replaces recurrence with \emph{self-attention}, in which every
token attends to every other token in the context window:
\begin{equation}
    \operatorname{Attention}(Q, K, V) =
    \operatorname{softmax}\!\left(\frac{QK^{\top}}{\sqrt{d_k}}\right)V,
\label{eq:attention}
\end{equation}
where $Q$, $K$, and $V$ are query, key, and value projections of the input
and $d_k$ is the key dimension. Stacked multi-head attention layers allow
the model to capture long-range dependencies --- in language and, as this
project exploits, in musical structure as well.

Brown et al.~\cite{brown20fewshot} demonstrated that sufficiently large
models exhibit \emph{in-context learning}: they can perform new tasks
specified entirely in the prompt, optionally guided by a small number of
examples (\emph{few-shot} prompting), without any parameter updates. This
\emph{zero-shot} or few-shot capability is the foundation of L2M. Modern
instruction-tuned models such as GPT-4 \cite{openai24gpt4} follow natural
language instructions and can emit structured output formats such as JSON on
request. Because their pre-training corpora include music theory texts, song
analyses, chord charts, and transcribed lyrics, such models implicitly
encode a substantial amount of musical knowledge --- knowledge that
music-specialized LLMs like ChatMusician \cite{yuan24chatmusician} make
explicit through further training, but which this work instead elicits
purely through prompting.

\section{Prompt Engineering}
\label{sec:bg-prompting}

Prompt engineering is the practice of designing model inputs that reliably
elicit a desired behavior. The techniques used in this work are:

\begin{itemize}
    \item \textbf{Role assignment.} A system message assigns the model a
    persona (e.g., ``expert music composer and emotional analyst''),
    conditioning it toward domain-appropriate output.
    \item \textbf{Output schema specification.} The prompt specifies an
    exact JSON structure for the response, allowing downstream code to parse
    and validate the output mechanically.
    \item \textbf{Few-shot examples.} A small number of worked input--output
    examples \cite{brown20fewshot} demonstrate the expected format and
    reasoning style; L2M uses three diverse examples per analysis prompt.
    \item \textbf{Hard constraints.} Explicit, non-negotiable requirements
    (e.g., ``generate EXACTLY $N$ notes'') that the output must satisfy.
    Decomposing a task into intermediate reasoning steps is likewise known
    to improve reliability on structured problems \cite{wei22cot}; L2M
    applies this principle architecturally by splitting analysis from
    generation.
    \item \textbf{Contextual carryover.} When input must be processed in
    chunks, appending a summary of previous output (here, the last three
    generated notes) to each subsequent prompt preserves continuity.
\end{itemize}

Even with these techniques, LLM output is stochastic and occasionally
malformed; every response in L2M is therefore validated against typed
schemas before use, and a deterministic fallback stands behind every LLM
call.

\section{Music, Emotion, and Tonality}
\label{sec:bg-emotion}

The link between musical structure and perceived emotion has been studied
empirically for nearly a century. Hevner's classic adjective-circle
experiments \cite{hevner36expression} established that listeners reliably
associate musical modes with affective qualities --- major modes with
happiness and serenity, minor modes with sadness and tension. Russell's
circumplex model of affect \cite{russell80circumplex} organizes emotion
along two dimensions, \emph{valence} (pleasant--unpleasant) and
\emph{arousal} (activated--deactivated), which map naturally onto musical
parameters: mode and consonance track valence, while tempo and dynamics
track arousal. Subsequent research consolidated these findings into robust
structure--emotion associations \cite{juslin01musicemotion,
gabrielsson10emotion}: fast tempi and ascending contours convey excitement
or joy; slow tempi, minor keys, and descending contours convey sadness;
moderate tempi with small intervals convey calm.

L2M operationalizes this literature directly: each of its six emotion
categories (\emph{happy, sad, hopeful, tense, calm, excited}) is mapped to a
key, a tempo range, and a melodic contour pattern
(\Autoref{chap:methodology}, \ref{tab:emotion-mapping}). Basic
tonal-music conventions constrain the output further: melodies remain
diatonic to the selected key, within a singable vocal range, and quantized
to standard note durations.

\section{Symbolic Music Representation}
\label{sec:bg-symbolic}

\subsection{Pitch, Duration, and Scientific Pitch Notation}

A monophonic melody is a sequence of notes, each defined by a pitch and a
duration. Pitches are written in \emph{scientific pitch notation} (SPN),
which combines a note letter, optional accidental, and octave number ---
e.g., C4 is middle C and A4 is the 440\,Hz tuning standard. Durations are
expressed in quarter-note (beat) units; L2M restricts them to the set
$\{0.25, 0.5, 1.0, 2.0, 4.0\}$, i.e., sixteenth note through whole note.
Loudness is encoded as MIDI \emph{velocity} (0--127).

\subsection{MIDI and MusicXML}

Two standard interchange formats are used for output. \textbf{MIDI}
(Musical Instrument Digital Interface) encodes music as timed note-on/note-off
events with pitch and velocity; it is universally supported by sequencers
and synthesizers. \textbf{MusicXML} \cite{good01musicxml} is an XML-based
notation interchange format that preserves score-level semantics --- key
signatures, time signatures, and note spellings --- making the output
directly editable in engraving software such as MuseScore or Finale.
Exporting both formats serves complementary use cases: playback and
production (MIDI) versus reading and editing by musicians (MusicXML).

\subsection{The music21 Toolkit}

music21 \cite{cuthbert10music21} is a Python toolkit for computer-aided
musicology that provides object models for notes, streams, key and time
signatures, and metronome marks, together with validated exporters for MIDI
and MusicXML. L2M builds its output stage on music21, which additionally
serves as an independent validity check: any generated file that music21
can re-parse without error is structurally well-formed.

\subsection{Audio Synthesis}

Symbolic output is optionally rendered to audio using
\emph{FluidSynth}, a SoundFont-based software synthesizer. A SoundFont
(\texttt{.sf2}) bundles sampled instrument timbres; FluidSynth renders a
MIDI file through these samples into PCM audio (WAV), which FFmpeg can
further encode to MP3. This completes the path from raw text to listenable
audio.
